%%% Основные сведения (на языке работы) / Main data (in the thesis language) %%%
% Эти поля печатаются на «главном» титульном листе: русском при thesislang=0 и
% английском при thesislang=1. Русский титульный лист английской работы берёт
% данные из блока «…Rus» ниже.
\ifnum\value{thesislang}=1
  %%%%%%%%%%%%%%%%%%%%%%%%%%%%%% English thesis %%%%%%%%%%%%%%%%%%%%%%%%%%%%%%%%%
  \newcommand{\thesisAuthorLastName}{\fixme{LastName}}
  \newcommand{\thesisAuthorOtherNames}{\fixme{First Middle}}
  \newcommand{\thesisAuthorInitials}{\fixme{F.\,M.}}
  \newcommand{\thesisTitle}            % English title
  {\fixme{Title of the dissertation}}
  \newcommand{\thesisSpecialtyTitle}   % English name of the specialty
  {\fixme{Name of the specialty}}
  \newcommand{\thesisDegree}           % English degree
  {\fixme{Doctor of Philosophy in \dots}}
  \newcommand{\thesisDegreeShort}      % English degree, short
  {\fixme{PhD}}
  \newcommand{\thesisCity}             % English city
  {\fixme{City}}
  \newcommand{\thesisOrganization}     % English organization
  {\fixme{Full name of the organization}}
  \newcommand{\thesisOrganizationShort}% English organization, short
  {\fixme{OrgShort}}
  \newcommand{\thesisInOrganization}   % English "The work was carried out at ..."
  {\fixme{the organization where the work was carried out}}
  \newcommand{\supervisorFio}          % Supervisor, name
  {\fixme{Supervisor Name}}
  \newcommand{\supervisorRegalia}      % Supervisor, regalia
  {\fixme{degree, title}}
  \newcommand{\supervisorFioShort}     % Supervisor, name, short
  {\fixme{F.\,M.~LastName}}
  \newcommand{\supervisorRegaliaShort} % Supervisor, regalia, short
  {\fixme{deg., tit.}}
\else
  %%%%%%%%%%%%%%%%%%%%%%%%%%%%%% Русская диссертация %%%%%%%%%%%%%%%%%%%%%%%%%%%%
  \newcommand{\thesisAuthorLastName}{\fixme{Фамилия}}
  \newcommand{\thesisAuthorOtherNames}{\fixme{Имя Отчество}}
  \newcommand{\thesisAuthorInitials}{\fixme{И.\,О.}}
  \newcommand{\thesisTitle}              % Диссертация, название
  {\fixme{Длинное название диссертационной работы, состоящее из~достаточно большого
  количества слов, совсем длинное длинное длинное длинное название, из~которого
  простому обывателю знакомы, в~лучшем случае, лишь отдельные слова}}
  \newcommand{\thesisSpecialtyTitle}     % Диссертация, специальность, название (название взято с сайта ВАК для примера)
  {\fixme{Технология обработки, хранения и~переработки злаковых, бобовых культур,
  крупяных продуктов, плодоовощной продукции и~виноградарства}}
  \newcommand{\thesisDegree}             % Диссертация, ученая степень
  {\fixme{кандидата физико-математических наук}}
  \newcommand{\thesisDegreeShort}        % Диссертация, ученая степень, краткая запись
  {\fixme{канд. физ.-мат. наук}}
  \newcommand{\thesisCity}               % Диссертация, город написания диссертации
  {\fixme{Город}}
  \newcommand{\thesisOrganization}       % Диссертация, организация
  {\fixme{Федеральное государственное автономное образовательное учреждение высшего
  образования <<Длинное название образовательного учреждения <<АББРЕВИАТУРА>>}}
  \newcommand{\thesisOrganizationShort}  % Диссертация, краткое название организации для доклада
  {\fixme{НазУчДисРаб}}
  \newcommand{\thesisInOrganization}     % Диссертация, организация в предложном падеже: Работа выполнена в ...
  {\fixme{учреждении с~длинным длинным длинным длинным названием, в~котором
  выполнялась данная диссертационная работа}}
  \newcommand{\supervisorFio}              % Научный руководитель, ФИО
  {\fixme{Фамилия Имя Отчество}}
  \newcommand{\supervisorRegalia}          % Научный руководитель, регалии
  {\fixme{уч. степень, уч. звание}}
  \newcommand{\supervisorFioShort}         % Научный руководитель, ФИО
  {\fixme{И.\,О.~Фамилия}}
  \newcommand{\supervisorRegaliaShort}     % Научный руководитель, регалии
  {\fixme{уч.~ст.,~уч.~зв.}}
\fi

\newcommand{\thesisAuthor}             % Диссертация, ФИО автора
{%
    \texorpdfstring{% \texorpdfstring takes two arguments and uses the first for (La)TeX and the second for pdf
        \thesisAuthorLastName~\thesisAuthorOtherNames% так будет отображаться на титульном листе или в тексте, где будет использоваться переменная
    }{%
        \thesisAuthorLastName, \thesisAuthorOtherNames% эта запись для свойств pdf-файла. В таком виде, если pdf будет обработан программами для сбора библиографических сведений, будет правильно представлена фамилия.
    }
}
\newcommand{\thesisAuthorShort}        % Диссертация, ФИО автора инициалами
{\thesisAuthorInitials~\thesisAuthorLastName}
%\newcommand{\thesisUdk}                % Диссертация, УДК
%{\fixme{xxx.xxx}}
\newcommand{\thesisSpecialtyNumber}    % Диссертация, специальность, номер (нейтрально)
{\fixme{XX.XX.XX}}
%% \newcommand{\thesisSpecialtyTwoNumber} % Диссертация, вторая специальность, номер
%% {\fixme{XX.XX.XX}}
%% \newcommand{\thesisSpecialtyTwoTitle}  % Диссертация, вторая специальность, название
%% {\fixme{Теория и~методика физического воспитания, спортивной тренировки,
%% оздоровительной и~адаптивной физической культуры}}
\newcommand{\thesisYear}               % Диссертация, год написания диссертации
{\the\year}

%% \newcommand{\supervisorDead}{}           % Рисовать рамку вокруг фамилии

%% \newcommand{\supervisorTwoDead}{}        % Рисовать рамку вокруг фамилии
%% \newcommand{\supervisorTwoFio}           % Второй научный руководитель, ФИО
%% {\fixme{Фамилия Имя Отчество}}
%% \newcommand{\supervisorTwoRegalia}       % Второй научный руководитель, регалии
%% {\fixme{уч. степень, уч. звание}}
%% \newcommand{\supervisorTwoFioShort}      % Второй научный руководитель, ФИО
%% {\fixme{И.\,О.~Фамилия}}
%% \newcommand{\supervisorTwoRegaliaShort}  % Второй научный руководитель, регалии
%% {\fixme{уч.~ст.,~уч.~зв.}}

\newcommand{\opponentOneFio}           % Оппонент 1, ФИО
{\fixme{Фамилия Имя Отчество}}
\newcommand{\opponentOneRegalia}       % Оппонент 1, регалии
{\fixme{доктор физико-математических наук, профессор}}
\newcommand{\opponentOneJobPlace}      % Оппонент 1, место работы
{\fixme{Не очень длинное название для места работы}}
\newcommand{\opponentOneJobPost}       % Оппонент 1, должность
{\fixme{старший научный сотрудник}}

\newcommand{\opponentTwoFio}           % Оппонент 2, ФИО
{\fixme{Фамилия Имя Отчество}}
\newcommand{\opponentTwoRegalia}       % Оппонент 2, регалии
{\fixme{кандидат физико-математических наук}}
\newcommand{\opponentTwoJobPlace}      % Оппонент 2, место работы
{\fixme{Основное место работы c длинным длинным длинным длинным названием}}
\newcommand{\opponentTwoJobPost}       % Оппонент 2, должность
{\fixme{старший научный сотрудник}}

%% \newcommand{\opponentThreeFio}         % Оппонент 3, ФИО
%% {\fixme{Фамилия Имя Отчество}}
%% \newcommand{\opponentThreeRegalia}     % Оппонент 3, регалии
%% {\fixme{кандидат физико-математических наук}}
%% \newcommand{\opponentThreeJobPlace}    % Оппонент 3, место работы
%% {\fixme{Основное место работы c длинным длинным длинным длинным названием}}
%% \newcommand{\opponentThreeJobPost}     % Оппонент 3, должность
%% {\fixme{старший научный сотрудник}}

\newcommand{\leadingOrganizationTitle} % Ведущая организация, дополнительные строки. Удалить, чтобы не отображать в автореферате
{\fixme{Федеральное государственное бюджетное образовательное учреждение высшего
профессионального образования с~длинным длинным длинным длинным названием}}

\newcommand{\defenseDate}              % Защита, дата
{\fixme{DD mmmmmmmm YYYY~г.~в~XX часов}}
\newcommand{\defenseDateEng}           % Защита, дата (англ.) --- только для thesislang=1
{\fixme{Month DD, YYYY~at~XX:XX}}
\newcommand{\defenseCouncilNumber}     % Защита, номер диссертационного совета
{\fixme{Д\,123.456.78}}
\newcommand{\defenseCouncilTitle}      % Защита, учреждение диссертационного совета
{\fixme{Название учреждения}}
\newcommand{\defenseCouncilAddress}    % Защита, адрес учреждение диссертационного совета
{\fixme{Адрес}}
\newcommand{\defenseCouncilAddressEng} % Защита, адрес (англ.) --- только для thesislang=1
{\fixme{Address in English}}
\newcommand{\defenseCouncilPhone}      % Телефон для справок
{\fixme{+7~(0000)~00-00-00}}

\newcommand{\defenseSecretaryFio}      % Секретарь диссертационного совета, ФИО
{\fixme{Фамилия Имя Отчество}}
\newcommand{\defenseSecretaryRegalia}  % Секретарь диссертационного совета, регалии
{\fixme{д-р~физ.-мат. наук}}            % Для сокращений есть ГОСТы, например: ГОСТ Р 7.0.12-2011 + http://base.garant.ru/179724/#block_30000

\newcommand{\synopsisLibrary}          % Автореферат, название библиотеки
{\fixme{Название библиотеки}}
\newcommand{\synopsisDate}             % Автореферат, дата рассылки
{\fixme{DD mmmmmmmm}\the\year~года}
\newcommand{\synopsisDateEng}          % Автореферат, дата рассылки (англ.) --- только для thesislang=1
{\fixme{DD mmmmmmmm}\the\year~year}

%%%%%%%%%%%%%%%%%%%%%%%%%%%%%%%%%%%%%%%%%%%%%%%%%%%%%%%%%%%%%%%%%%%%%%%%%%%%%%%%
%%% Обязательные русские сведения для английской диссертации (thesislang=1) %%%
%%% Russian data for the mandatory Russian title page & abstract (EN thesis) %
%%%%%%%%%%%%%%%%%%%%%%%%%%%%%%%%%%%%%%%%%%%%%%%%%%%%%%%%%%%%%%%%%%%%%%%%%%%%%%%%
% Русская диссертация (thesislang=0): эти команды НЕ используются --- заполнять
%   нужно только основные сведения выше.
% Английская диссертация (thesislang=1): основные сведения выше заполняются
%   по-английски, а перечисленные ниже --- по-русски (для обязательного русского
%   титульного листа и аннотации).
\newcommand{\thesisAuthorRus}          % ФИО автора по-русски (полностью)
{\fixme{Фамилия Имя Отчество}}
\newcommand{\thesisTitleRus}           % Название диссертации по-русски
{\fixme{Русское название диссертационной работы}}
\newcommand{\thesisSpecialtyTitleRus}  % Название специальности по-русски
{\fixme{Русское название специальности}}
\newcommand{\thesisDegreeRus}          % Учёная степень по-русски (род. падеж)
{\fixme{кандидата физико-математических наук}}
\newcommand{\thesisCityRus}            % Город по-русски
{\fixme{Город}}
\newcommand{\thesisOrganizationRus}    % Организация по-русски (для шапки титула)
{\fixme{Полное название организации по-русски}}
\newcommand{\supervisorFioRus}         % Научный руководитель, ФИО по-русски
{\fixme{Фамилия Имя Отчество}}
\newcommand{\supervisorRegaliaRus}     % Научный руководитель, регалии по-русски
{\fixme{уч. степень, уч. звание}}

% To avoid conflict with beamer class use \providecommand
\providecommand{\keywords}%            % Ключевые слова для метаданных PDF диссертации и автореферата
{}
